\pgfplotsset{compat=1.18}
% Cores para o estilo científico
\definecolor{myblue}{RGB}{0, 114, 189}     % Azul de um paleta comum em ciência
\definecolor{myred}{RGB}{217, 83, 25}      % Vermelho alaranjado
\definecolor{mygreen}{RGB}{93, 169, 58}    % Verde escuro
\definecolor{mypurple}{RGB}{128, 0, 128}   % Cor para Regressão-2
\definecolor{myorange}{RGB}{255, 140, 0} % Nova cor para Regressão-3 (exemplo: laranja)

\begin{filecontents*}{dados.csv}
n, k, densidade, m_prime
20, 7.7251, 0.1421, 41.8
40, 6.4208, 0.0641, 70.8
60, 4.6691, 0.0401, 95.0
80, 3.7304, 0.0294, 119.8
100, 3.1524, 0.0230, 144.0
120, 3.1734, 0.0190, 168.8
140, 2.7514, 0.0161, 193.0
160, 2.9443, 0.0141, 217.8
180, 2.5279, 0.0124, 242.0
200, 2.5287, 0.0112, 266.8
250, 2.3600, 0.0088, 321.0
300, 2.1822, 0.0073, 380.0
350, 2.1600, 0.0062, 434.8
400, 2.0859, 0.0054, 489.0
450, 2.0559, 0.0048, 548.0
500, 2.0122, 0.0043, 602.8
600, 1.6702, 0.0035, 711.8
700, 1.6548, 0.0030, 820.8
800, 1.6506, 0.0026, 929.8
900, 1.6330, 0.0023, 1034.0
1000, 1.6196, 0.0021, 1143.0
\end{filecontents*}

\begin{figure}[htbp]
    \centering
    \resizebox{1\textwidth}{!}{
        \begin{tikzpicture}
            \begin{axis}[
                xlabel={Vértices ($n$)},
                ylabel={$K$},
                title={Relação entre Vértices e $K$},
                xmin=0, xmax=1100,
                ymin=1.5, ymax=8.0,
                xtick={0,200,400,600,800,1000},
                ytick={1.5,2.0,2.5,3.0,3.5,4.0,4.5,5.0,5.5,6.0,6.5,7.0,7.5,8.0},
                % Ajustes para estilo científico
                legend pos=north east,
                ymajorgrids=true,
                grid style={dashed, gray!50}, % Grade mais suave
                axis lines=box, % Desenha as bordas do box
                tick style={color=black}, % Cor dos ticks
                label style={font=\bfseries}, % Rótulos em negrito
                title style={font=\bfseries}, % Título em negrito
                % Ajustando a largura e altura
                width=12cm,
                height=8cm,
            ]
                % Adicionando curva suavizada (para os dados de K Exato)
                \addplot[
                    smooth,
                    thick,
                    color=myred, % Cor vermelha para a curva suavizada (K Exato)
                    mark=none
                ] table [x=n, y=alfa, col sep=comma] {dados.csv};

                % Adicionando a reta de regressão 1
                \addplot[
                    color=mygreen, % Cor verde para a regressão 1
                    thick,
                    dash dot, % Estilo da linha (alternando traço e ponto para distinção)
                    mark=none
                ] table [
                    x=n,
                    y expr={32.5649 * \thisrow{densidade} + 1.9948},
                    col sep=comma
                ] {dados.csv};

                % Adicionando a reta de regressão 2
                \addplot[
                    color=mypurple, % Nova cor para a Regressão-2 (exemplo: roxo)
                    thick,
                    dotted, % Estilo da linha (apenas pontos para distinção)
                    mark=none
                ] coordinates {
                    (0, { -0.0036 * 0 + 4.1587 }) % Ponto inicial (n=0)
                    (1000, { -0.0036 * 1000 + 4.1587 }) % Ponto final (n=1000)
                };

                % Adicionando a reta de regressão 3
                \addplot[
                    color=myorange, % Nova cor para a Regressão-3 (exemplo: laranja)
                    thick,
                    densely dashed, % Outro estilo de linha para distinção
                    mark=none
                ] table [
                    x=n,
                    y expr={-0.0098 * \thisrow{n} + 0.0055 * \thisrow{m_prime} + 3.9640},
                    col sep=comma
                ] {dados.csv};

            \legend{K\_Exato, Regressão\_1, Regressão\_2, Regressão\_3}
            \end{axis}
        \end{tikzpicture}
    }
    \caption{Relação entre o valor de $K$ (número de vértices branch) e o número de vértices totais ($n$) em grafos. A curva suavizada representa a tendência dos dados exatos, e as linhas de regressão mostram modelos lineares de $K$ em função da densidade ou do número de vértices.}
    \label{fig:grafic_alfa}
\end{figure}